\section{Линейные отображения и их координатное представление}

\begin{definition}[Линейное отображение]
Отображение $T \colon V \to W$ между вещественными векторными пространствами называется \textbf{линейным отображением}
(линейным оператором / гомоморфизмом), если для любых $u, v \in V$ и любого $\alpha \in \mathbb{R}$ выполнены:
\begin{enumerate}
    \item Аддитивность: $T(u + v) = T(u) + T(v)$;
    \item Однородность: $T(\alpha u) = \alpha T(u)$.
\end{enumerate}
Эквивалентно сохранению произвольных линейных комбинаций:
\[ T\left(\sum_{i=1}^k c_i v_i\right) = \sum_{i=1}^k c_i T(v_i). \]
\end{definition}

\begin{definition}[Ядро и образ]
Пусть $T \colon V \to W$ — линейное отображение.
\begin{itemize}
    \item \textbf{Ядром} (\textbf{kernel / nullspace}) отображения $T$ называется множество векторов из $V$,
    переходящих в нулевой вектор пространства $W$: \[ \ker(T) = \{ v \in V \mid T(v) = 0_W \} \le V \].
    \item \textbf{Образом} (\textbf{image / range}) отображения $T$ называется подмножество векторов из $W$, которые
    достигаются отображением $T$:
    \[ \im(T) = \{ w \in W \mid \exists v \in V \colon T(v) = w \} \le W. \]
\end{itemize}
Отображение $T$ \textbf{инъективно} тогда и только тогда, когда $\ker(T) = \{0_V\}$.\\ Отображение $T$
\textbf{сюръективно} тогда и только тогда, когда $\im(T) = W$.
\end{definition}

\section{Матрица линейного отображения: почему столбцы суть образы базисных векторов?}

Пусть в пространстве $V$ зафиксирован базис $\mathcal{B}=(v_1,\dots, v_n)$, а в пространстве $W$ зафиксирован базис
$\mathcal{C} = (w_1, \dots, w_m)$.

Возьмём произвольный вектор $x \in V$.
Разложим его по базису $\mathcal{B}$:
\[ x = x_1 v_1 + x_2 v_2 + \dots + x_n v_n = \sum_{j=1}^n x_j v_j. \]
Применим линейное отображение $T$ к вектору $x$, используя свойство линейности:
\[ T(x) = T\left(\sum_{j=1}^n x_j v_j\right) = \sum_{j=1}^n x_j T(v_j). \]
Отсюда виден ключевой фундаментальный факт: \textbf{действие линейного отображения на любой вектор полностью и
однозначно определяется тем, куда переходят базисные векторы} $T(v_1), \dots, T(v_n) \in W$.

Теперь разложим каждый вектор $T(v_j) \in W$ по базису $\mathcal{C}$:
\[ T(v_j) = \sum_{i=1}^m a_{ij} w_i = a_{1j} w_1 + a_{2j} w_2 + \dots + a_{mj} w_m. \]
Координатный столбец вектора $T(v_j)$ в базисе $\mathcal{C}$ равен:
\[ [T(v_j)]_\mathcal{C} = \begin{pmatrix}
                              a_{1j} \\ a_{2j} \\ \vdots \\ a_{mj}
\end{pmatrix}. \]
Подставим это разложение в выражение для $T(x)$:
\[ T(x)=\sum_{j=1}^n x_j\left(\sum_{i=1}^m a_{ij} w_i\right)=\sum_{i=1}^m \left( \sum_{j=1}^n a_{ij} x_j\right) w_i.
\]
Следовательно , $i$-я координата вектора $y = T(x)$ в базисе $\mathcal{C}$ равна $y_i = \sum_{j=1}^n a_{ij} x_j$.
В матричной форме это записывается как:
\[ \begin{pmatrix}
       y_1 \\ \vdots \\ y_m
\end{pmatrix} = \begin{pmatrix}
                    a_{11} & a_{12} & \dots & a_{1n} \\ a_{21} & a_{22} & \dots & a_{2n} \\ \vdots
                    & \vdots & \ddots & \vdots \\ a_{m1} & a_{m2} & \dots & a_{mn}
\end{pmatrix} \begin{pmatrix}
                  x_1 \\ \vdots \\ x_n
\end{pmatrix} \iff [T(x)]_\mathcal{C} = [T]_{\mathcal{C} \leftarrow
\mathcal{B}} [x]_\mathcal{B}. \]

\begin{theorem}[Структура матрицы линейного оператора]
Матрица линейного оператора $[T]_{\mathcal{C} \leftarrow \mathcal{B}} \in \mathbb{R}^{m \times n}$ состоит из столбцов,
каждый $j$-й столбец которых является координатным представлением образа $j$-го базисного вектора $v_j$ в базисе
$\mathcal{C}$:
\[ [T]_{\mathcal{C} \leftarrow \mathcal{B}} = \begin{pmatrix}
                                                  | & | & & | \\ [T(v_1)]_\mathcal{C} & [T(v_2)]_\mathcal{C}
                                                  & \dots & [T(v_n)]_\mathcal{C} \\
                                                  | & | & & |
\end{pmatrix}. \]
\end{theorem}

\section{Концептуальный смысл теоремы о ранге и дефекте (Rank-Nullity)}

\begin{theorem}[Rank-Nullity Theorem]
Пусть $V$ и $W$ — конечномерные векторные пространства, и $T \colon V \to W$ — линейное отображение.
Тогда:
\[ \dim V = \dim(\ker T) + \dim(\im T). \]
В терминах ранга и дефекта (дефект $\nullity(T) \coloneqq \dim(\ker T)$, ранг $\rank(T) \coloneqq \dim(\im T)$):
\[ \dim V = \rank(T) + \nullity(T). \]
Для матрицы $A \in \mathbb{R}^{m \times n}$ с $n$ столбцами:
\[ n = \rank(A) + \dim(\Null A). \]
\end{theorem}

\begin{proof}[Идея доказательства]
Пусть $\dim V = n$, и пусть $\dim(\ker T) = k$.
Выберем базис ядра $\ker T$: пусть это векторы $(u_1, u_2, \dots, u_k)$.
Поскольку эти векторы линейно независимы в $V$, по теореме о дополнении базиса их можно дополнить векторами $(v_1,
\dots, v_{n-k})$ до полного базиса всего пространства $V$:
\[ \mathcal{B} = (u_1, \dots, u_k, v_1, \dots, v_{n-k}). \]
Рассмотрим образы дополнительных векторов в $W$: $w_j \coloneqq T(v_j)$ для $j = 1, \dots, n-k$.
Покажем, что система $(w_1, \dots, w_{n-k})$ образует базис подпространства $\im T$:
\begin{enumerate}
    \item \textbf{Порождаемость:} Любой вектор $x \in V$ раскладывается как $x = \sum_{i=1}^k \alpha_i u_i +
    \sum_{j=1}^{n-k} \beta_j v_j$.
    Применяя $T$ и учитывая, что $T(u_i) = 0$:
    \[ T(x) = \sum_{i=1}^k \alpha_i T(u_i) + \sum_{j=1}^{n-k} \beta_j T(v_j) = 0 + \sum_{j=1}^{n-k} \beta_j w_j. \]
    Следовательно, векторы $w_1, \dots, w_{n-k}$ порождают $\im T$.
    \item \textbf{Линейная независимость:} Пусть $\sum_{j=1}^{n-k} c_j w_j = 0$.
    Тогда $T\left(\sum_{j=1}^{n-k} c_j
    v_j\right) = 0$, что означает, что вектор $v \coloneqq \sum_{j=1}^{n-k} c_j v_j$ лежит в ядре $\ker T$.
    Но тогда он должен выражаться через базис ядра $(u_1, \dots, u_k)$:
    \[ \sum_{j=1}^{n-k} c_j v_j = \sum_{i=1}^k d_i u_i \iff \sum_{j=1}^{n-k} c_j v_j - \sum_{i=1}^k d_i u_i = 0. \]
    Так как векторы $(u_1,\dots, u_k,v_1,\dots, v_{n-k})$ образуют базис $V$, они линейно независимы, следовательно, все
    коэффициенты $c_j = 0$ и $d_i = 0$.
\end{enumerate}
Итак, $(w_1, \dots, w_{n-k})$ — базис $\im T$, а значит, $\dim(\im T) = n - k$.
Складывая: $\dim(\ker T) + \dim(\im T) = k + (n - k) = n = \dim V$.
\end{proof}

\begin{remark}[Концептуальная геометрическая интерпретация]
Теорема о ранге и дефекте выражает \textbf{закон сохранения размерности}: Исходное пространство размерности $n$ при
линейном отображении делится на две непересекающиеся составляющие:
\begin{itemize}
    \item Направления, которые отображение \textbf{схлопывает в ноль} (ядро $\ker T$ размерности $\nullity T$).
    \item Направления, которые \textbf{переживают трансформацию} и разворачиваются в образе
    ($\im T$ размерности $\rank T$).
\end{itemize}
Ни одно измерение не может исчезнуть бесследно: либо степень свободы уничтожается отображением (уходит в дефект), либо
проявляется на выходе (даёт вклад в ранг).
\end{remark}

\section{Четыре фундаментальных подпространства матрицы}

Для матрицы $A \in \mathbb{R}^{m \times n}$ канонически определены 4 подпространства (теорема Странга):
\begin{enumerate}
    \item \textbf{Пространство столбцов} (\textbf{Column space}): $\Col(A) = \im(A) \le \mathbb{R}^m$,
    размерность $\rank(A) = r$;
    \item \textbf{Пространство строк} (\textbf{Row space}): $\Row(A) = \Col(A^\top) \le \mathbb{R}^n$, размерность
    $\rank(A^\top) = \rank(A) = r$;
    \item \textbf{Нуль-пространство (ядро)} (\textbf{Nullspace}): $\Null(A) = \ker(A) \le \mathbb{R}^n$,
    размерность $n - r$;
    \item \textbf{Левое нуль-пространство}: $\Null(A^\top) = \ker(A^\top) \le \mathbb{R}^m$, размерность $m - r$.
\end{enumerate}

\begin{theorem}[Ортогональные дополнения фундаментальных подпространств]
В пространствах $\mathbb{R}^n$ и $\mathbb{R}^m$ фундаментальные подпространства взаимно ортогональны и образуют полные
прямые суммы:
\begin{gather*}
    \Null(A) = (\Row(A))^\perp \quad \text{в} \quad \mathbb{R}^n \implies \mathbb{R}^n = \Row(A) \oplus \Null(A),\\
    \Null(A^\top) = (\Col(A))^\perp \quad \text{в} \quad \mathbb{R}^m \implies \mathbb{R}^m = \Col(A) \oplus
    \Null(A^\top).\\
\end{gather*}
\end{theorem}

\section{Диагностика невырожденности: определитель, ранг и обратимость}

Для квадратной матрицы $A\in\mathbb{R}^{n\times n}$ понятия определителя, ранга и обратимости представляют собой три
взаимодополняющих взгляда на одно и то же явление — \textbf{невырожденность оператора}:

\begin{theorem}[Теорема об эквивалентности условий невырожденности]
Для квадратной матрицы $A \in \mathbb{R}^{n \times n}$ следующие утверждения строго эквивалентны:
\begin{enumerate}
    \item Матрица $A$ обратима (существует $A^{-1}$ такая, что $A A^{-1} = A^{-1} A = I_n$);
    \item Определитель отличен от нуля: $\det(A) \ne 0$;
    \item Матрица $A$ имеет полный ранг: $\rank(A) = n$;
    \item Ядро матрицы тривиально: $\Null(A) = \{0\}$;
    \item Столбцы матрицы $A$ линейно независимы и образуют базис $\mathbb{R}^n$;
    \item Строки матрицы $A$ линейно независимы и образуют базис $\mathbb{R}^n$;
    \item Система линейных уравнений $Ax = 0$ имеет исключительно тривиальное решение $x = 0$;
    \item Для любого вектора $b \in \mathbb{R}^n$ система $Ax = b$ имеет единственное решение $x = A^{-1} b$;
    \item Число 0 не является собственным значением матрицы $A$ ($\lambda_i \ne 0$ для всех $i$);
    \item Все сингулярные числа матрицы $A$ строго положительны ($\sigma_n > 0$).
\end{enumerate}
\end{theorem}

\textbf{Геометрическая интерпретация проблемы:}
\begin{itemize}
    \item $\det(A) = 0$ означает, что оператор сжимает $n$-мерный объём единичного гиперкуба в нуль
    (отображение сплющивает пространство в подпространство меньшей размерности).
    \item $\rank(A) < n$ указывает размерность полученного сплющенного образа.
    \item Несуществование $A^{-1}$ констатирует невозможность обратить этот процесс вспять
    (информация о направлении, попавшем в ядро, безвозвратно утеряна).
\end{itemize}

\section{Системы линейных уравнений и элементарные операции}

Система $m$ линейных уравнений с $n$ неизвестными записывается в виде:
\[ Ax = b, \quad A \in \mathbb{R}^{m \times n}, \; x \in \mathbb{R}^n, \; b \in \mathbb{R}^m. \]
\begin{theorem}[Теорема Кронекера-Капелли]
Система линейных уравнений $Ax = b$ совместна (имеет хотя бы одно решение) тогда и только тогда, когда ранг основной
матрицы равен рангу расширенной матрицы:
\[ \rank(A) = \rank(A \mid b). \]
При этом:
\begin{enumerate}
    \item Если $\rank(A) = \rank(A \mid b) = n$ (число неизвестных), система имеет \textbf{ровно одно решение}.
    \item Если $\rank(A) = \rank(A \mid b)=r < n$, система имеет \textbf{бесконечно много решений}, зависящих от $n - r$
    свободных параметров.
    \item Если $\rank(A) < \rank(A \mid b)$, система \textbf{несовместна} (решений нет).
\end{enumerate}
\end{theorem}

\subsection{Почему элементарные преобразования строк не изменяют решений?}
\begin{definition}[Элементарные строчные операции]
Существуют три типа элементарных преобразований строк матрицы:
\begin{enumerate}
    \item Перестановка двух строк местами ($R_i \leftrightarrow R_j$);
    \item Умножение строки на ненулевой скаляр ($R_i \to c R_i$, $c \ne 0$);
    \item Прибавление к одной строке другой строки, умноженной на скаляр ($R_i \to R_i + c R_j$).
\end{enumerate}
\end{definition}

\begin{theorem}[О сохранении множества решений]
Элементарные преобразования строк расширенной матрицы $(A \mid b)$ не изменяют множества решений системы $Ax = b$.
\end{theorem}
\begin{proof}
    Любое элементарное преобразование строк эквивалентно умножению расширенной матрицы слева на соответствующую
    \textbf{элементарную матрицу} $E \in \mathbb{R}^{m \times m}$:
    \[ E(A \mid b) = (EA \mid Eb). \]
    Фундаментальное свойство элементарных матриц заключается в том, что \textbf{все они строго обратимы},
    причём матрица $E^{-1}$ сама является элементарной матрицей того же типа:
    \begin{itemize}
        \item Матрица перестановки строк $P_{ij}$ совпадает со своей обратной: $P_{ij}^{-1} = P_{ij}$;
        \item Обратная к умножению строки на $c \ne 0$ есть умножение на $\frac{1}{c}$;
        \item Обратная к прибавлению строки $j$, умноженной на $c$, есть вычитание строки $j$, умноженной на $c$.
    \end{itemize}
    Пусть $x$ — решение исходной системы, то есть $Ax = b$.
    Умножив обе части равенства слева на $E$, получаем $E(Ax) = Eb \iff (EA)x = Eb$.
    Следовательно, $x$ является решением преобразованной системы.
    
    Обратно, пусть $x$ — решение преобразованной системы, то есть $(EA)x = Eb$.
    Поскольку существует обратная матрица $E^{-1}$, умножим равенство слева на $E^{-1}$:
    \[ E^{-1}((EA)x) = E^{-1}(Eb) \iff (E^{-1}E)Ax = (E^{-1}E)b \iff Ax = b. \]
    Таким образом, $x$ является решением исходной системы.
    Множества решений совпадают.
\end{proof}

\section{LU-разложение и его вычислительное преимущество}

\begin{definition}[LU-decomposition]
\textbf{$LU$-разложением} квадратной матрицы $A \in \mathbb{R}^{n \times n}$ называется представление вида:
\[ A = L U, \]
где $L \in \mathbb{R}^{n \times n}$ — нижнетреугольная матрица (с единицами на главной диагонали в случае разложения
Дулиттла), а $U \in \mathbb{R}^{n \times n}$ — верхнетреугольная матрица.
Если в процессе исключения Гаусса требуются перестановки строк (выбор главного элемента для численной устойчивости),
используется разложение с перестановкой:
\[ P A = L U, \]
где $P$ — матрица перестановки.
\end{definition}

\subsection{Вычислительное преимущество LU-разложения}
Пусть требуется решить систему $Ax = b$.
Подставив $A = LU$, получаем:
\[ L (U x) = b. \]
Решение распадается на два последовательных простых треугольных этапа:
\begin{enumerate}
    \item \textbf{Прямая подстановка (Forward substitution):} решаем систему с нижнетреугольной матрицей $L y = b$.
    Требует $\approx \frac{n^2}{2}$ операций.
    \item \textbf{Обратная подстановка (Back substitution):} решаем систему с верхнетреугольной матрицей $U x = y$.
    Требует $\approx \frac{n^2}{2}$ операций.
\end{enumerate}

\textbf{Сравнение сложности:}
\begin{itemize}
    \item Прямой метод Гаусса «с нуля» требует $\mathcal{O}\left(\frac{2}{3} n^3\right)$ арифметических операций
    для каждого нового вектора правой части $b$.
    \item Факторизация $A=LU$ вычисляется \textbf{один раз} со сложностью $\mathcal{O}\left(\frac{2}{3} n^3\right)$.
    \item После этого для любого нового вектора $b$ нахождение решения через $Ly = b$ и $Ux = y$ требует всего лишь
    $\mathcal{O}(n^2)$ операций!
\end{itemize}
В реальных прикладных задачах (симуляция динамических систем, нелинейные итерации метода Ньютона, обработка
сигналов), где матрица $A$ неизменна, а правая часть $b$ меняется тысячи раз за секунду, использование
$LU$-разложения ускоряет расчёты в тысячи раз.

\section{Типичные ошибки размерностей при работе с матрицами}

\begin{enumerate}
    \item \textbf{Несогласованность внутренних размерностей при умножении:} Матричное произведение $A B$
    определено тогда и только тогда, когда число столбцов левой матрицы совпадает с числом строк правой матрицы:
    \[ A_{(m \times k)} \cdot B_{(k \times n)} = C_{(m \times n)}. \]
    Попытка умножить $A_{(m \times k)}$ на $B_{(p \times n)}$ при $k \ne p$ является грубейшей синтаксической и
    содержательной ошибкой.
    \item \textbf{Путаница между столбцами ($n \times 1$) и строками ($1 \times n$):} В математике по умолчанию
    вектор $x \in \mathbb{R}^n$ является столбцом ($n \times 1$).
    \begin{itemize}
        \item $x^\top y$ — \textbf{внутреннее (скалярное) произведение}: матрица
        $(1 \times n) \times (n \times 1) = (1 \times 1) \implies \text{число}$;
        \item $x y^\top$ — \textbf{внешнее произведение (outer product)}: матрица
        $(n \times 1) \times (1 \times m) = (n \times m) \implies \text{матрица ранга 1}$.
    \end{itemize}
    \item \textbf{Ошибки транспонирования произведения:}
    \[ (A B)^\top = B^\top A^\top \quad (\text{порядок сомножителей обязательно обращается!}). \]
    Написание $(AB)^\top = A^\top B^\top$ в общем случае приводит даже к несовпадению размерностей.
    \item \textbf{Ложное предположение о коммутативности:} В общем случае $A B \ne B A$.
    Более того, если $A B$ определено, произведение $B A$ может даже не существовать из-за несовпадения
    размерностей.
\end{enumerate}

\newpage

% =========================================================================
% ГЛАВА 3
% =========================================================================